\chapter{Derivaciones Matemáticas}
\label{ap:derivaciones}

Este apéndice reúne las derivaciones matemáticas necesarias para interpretar los métodos GSP y TRSS usados en la tesis. El objetivo no es repetir el Capítulo~\ref{ch:marco_teorico}, sino explicitar las formas algebraicas que justifican la implementación computacional.

\section{Solución de Tikhonov en Grafos}

Sea $\mathbf{y}\in\mathbb{R}^{N}$ la señal observada en un instante y sea $\mathbf{M}=\operatorname{diag}(m_1,\dots,m_N)$ una máscara diagonal con $m_i=1$ para canales observados y $m_i=0$ para canales faltantes. La regularización de Tikhonov en grafos estima
\begin{equation}
\hat{\mathbf{x}}=\arg\min_{\mathbf{x}}\; \|\mathbf{M}(\mathbf{x}-\mathbf{y})\|_2^2 + \alpha\,\mathbf{x}^{\top}\mathbf{L}\mathbf{x},
\end{equation}
donde $\mathbf{L}$ es el Laplaciano del grafo y $\alpha>0$ controla el peso de suavidad espacial. Derivando respecto de $\mathbf{x}$ e igualando a cero:
\begin{equation}
2\mathbf{M}^{\top}\mathbf{M}(\mathbf{x}-\mathbf{y}) + 2\alpha\mathbf{L}\mathbf{x}=0.
\end{equation}
Como $\mathbf{M}$ es diagonal binaria, $\mathbf{M}^{\top}\mathbf{M}=\mathbf{M}$. Por tanto,
\begin{equation}
(\mathbf{M}+\alpha\mathbf{L})\hat{\mathbf{x}}=\mathbf{M}\mathbf{y},
\qquad
\hat{\mathbf{x}}=(\mathbf{M}+\alpha\mathbf{L})^{-1}\mathbf{M}\mathbf{y}.
\end{equation}
En la práctica no es necesario formar la inversa explícita: basta resolver el sistema lineal. La presencia de $\alpha\mathbf{L}$ acopla los canales faltantes con los canales observados según la topología del grafo.

\section{Solución de TRSS en Forma Matricial}

TRSS extiende la regularización anterior a una ventana temporal $\mathbf{X}\in\mathbb{R}^{N\times T}$. La formulación usada en esta tesis es
\begin{equation}
\hat{\mathbf{X}}=\arg\min_{\mathbf{X}}\; \|\mathbf{M}\odot(\mathbf{X}-\mathbf{Y})\|_F^2 + \alpha\operatorname{tr}(\mathbf{X}^{\top}\mathbf{L}\mathbf{X}) + \beta\|\mathbf{X}\mathbf{D}_t\|_F^2,
\end{equation}
donde $\mathbf{D}_t$ es el operador de diferencias finitas temporales. Al vectorizar $\mathbf{X}$ como $\tilde{\mathbf{x}}=\operatorname{vec}(\mathbf{X})$, el término espacial se expresa mediante $\mathbf{I}_T\otimes\mathbf{L}$ y el término temporal mediante $\mathbf{D}_t\mathbf{D}_t^{\top}\otimes\mathbf{I}_N$. El sistema lineal resultante es
\begin{equation}
\left(\tilde{\mathbf{M}} + \alpha(\mathbf{I}_T\otimes\mathbf{L}) + \beta(\mathbf{D}_t\mathbf{D}_t^{\top}\otimes\mathbf{I}_N)\right)\hat{\tilde{\mathbf{x}}}=\tilde{\mathbf{M}}\tilde{\mathbf{y}}.
\end{equation}
Esta forma muestra por qué TRSS es más costoso que una interpolación puramente espacial: el sistema acopla canales y tiempo. También muestra su ventaja conceptual: la reconstrucción no se decide muestra a muestra, sino considerando continuidad temporal de la ventana completa.

\section{Operador de Diferencias Temporales}

El operador de diferencias finitas hacia adelante $\mathbf{D}_t\in\mathbb{R}^{T\times(T-1)}$ contiene $-1$ y $+1$ en posiciones consecutivas. Para una fila temporal $\mathbf{x}_{i,:}$, el producto $\mathbf{x}_{i,:}\mathbf{D}_t$ calcula diferencias $x_{i,t+1}-x_{i,t}$. Por ello,
\begin{equation}
\|\mathbf{X}\mathbf{D}_t\|_F^2 = \sum_{i=1}^{N}\sum_{t=1}^{T-1}(X_{i,t+1}-X_{i,t})^2.
\end{equation}
El término penaliza saltos abruptos y favorece reconstrucciones temporalmente suaves. No agrega información externa; restringe el conjunto de soluciones compatibles con los canales observados.

\section{Propiedades Espectrales del Laplaciano}

Para un grafo no dirigido con pesos no negativos, el Laplaciano $\mathbf{L}=\mathbf{D}-\mathbf{W}$ es semidefinido positivo. Su forma cuadrática cumple
\begin{equation}
\mathbf{x}^{\top}\mathbf{L}\mathbf{x}=\frac{1}{2}\sum_{i,j}W_{ij}(x_i-x_j)^2\geq 0.
\end{equation}
El primer autovalor es cero y su autovector asociado es constante en cada componente conexa. La conectividad algebraica $\lambda_2$ aumenta cuando el grafo está mejor conectado. En interpolación, un grafo demasiado desconectado limita la propagación de información, mientras que un grafo demasiado denso puede suavizar en exceso. Esta tensión justifica la evaluación empírica de varios grafos en la selección preliminar.

\section{Implicancias Computacionales}

Las derivaciones anteriores explican la diferencia práctica entre Tikhonov espacial, TRSS y MNE-Python. En Tikhonov espacial, el sistema se resuelve de forma independiente para cada muestra temporal, por lo que el costo principal depende de $N$ y de la estructura de $\mathbf{L}$. En TRSS, en cambio, el operador vectorizado acopla $N$ canales y $T$ muestras; por ello, aun cuando las matrices sean dispersas, el resolvedor opera sobre un espacio de mayor dimensión efectiva.

Esta diferencia no invalida TRSS, pero sí condiciona su uso. Si una ventana temporal es corta, la pérdida espacial es severa y el análisis se realiza fuera de línea, el costo adicional puede ser aceptable. Si se procesan muchos sujetos, muchas ventanas o flujos de baja latencia, conviene reutilizar factorizaciones, explotar dispersidad y limitar la búsqueda de hiperparámetros. La frontera práctica descrita en el Capítulo~\ref{ch:resultados} se deriva precisamente de esta combinación entre forma matemática, calidad de reconstrucción y costo computacional.
